% !PRINTED 552
\weberpage{552}
% running head: 552 Siebzehnter Abschnitt. \S.\,152.
%<<CONT
theilbare Zahl sich als Product von idealen Zahlen darstellen
lasse. Alle diese Fragen werden beantwortet durch die Vermittelung
 der Functionaltheorie (oder der Idealtheorie), wenn man
der idealen Zahl $(\alpha, \beta)$ das Functional $\varphi$ zuordnet, und können
wohl kaum auf einfachere Weise entschieden werden. Die Einführung
 des Systemes der idealen Zahlen bietet dann auch keine
wesentlichen Vortheile mehr, und man kann ebenso gut die
Functionale selbst, oder die Ideale als Elemente der neu einzuführenden
 Mannigfaltigkeit betrachten.

\begin{center}
\S.\,152.

\gesperrt{Aequivalenz.}
\end{center}

Wir nennen jetzt zwei Functionale, die sich nur durch einen
Einheitsfactor unterscheiden, auch wenn sie gebrochen sind,
\gesperrt{associirt}, und stellen folgende Definition auf:

1. \gesperrt{Zwei ganze oder gebrochene Functionale $\varphi, \psi$
im Körper $\Omega$ heissen äquivalent, wenn ihr
Quotient $\varphi : \psi$ mit einer Zahl associirt ist.}

Es heissen also die beiden Functionale $\varphi$ und $\psi$ äquivalent,
wenn eine Einheit $\varepsilon$ und eine Zahl $\alpha$ in $\Omega$ existiren, so dass

\weq{1}{\frac{\varphi}{\psi} = \alpha\varepsilon}

ist. Ist $\varphi$ äquivalent mit $\psi$ und mit $\psi_1$, so folgt aus (1)

\[
\frac{\varphi}{\psi} = \alpha\varepsilon, \qquad \frac{\varphi}{\psi_1} = \alpha_1\varepsilon_1,
\]

folglich

\[
\frac{\psi}{\psi_1} = \frac{\alpha_1}{\alpha} \frac{\varepsilon_1}{\varepsilon},
\]

und da $\alpha_1 : \alpha$ eine Zahl, $\varepsilon_1 : \varepsilon$ eine Einheit ist, so folgt der
erste Satz:

2. \gesperrt{Zwei Functionale, die mit einem dritten äquivalent
 sind, sind auch unter einander äquivalent
.}

Theilt man hiernach alle Functionale des Körpers $\Omega$ in
\gesperrt{Classen} ein, indem man zwei Functionale in dieselbe oder in
verschiedene Classen wirft, je nachdem sie äquivalent sind oder
nicht, so ergiebt sich, dass zwei dieser Classen, die ein einziges
%CONT>>
